%! The Normal Distribution and z-Scores — Unit 1, Lesson 1.7 — AP Practice (Answer Key)
% EFFL, AP Statistics variant: AP-format practice, assigned but NOT scored — the cover's score
% column reads NA for this component. The graded practice is the `homework` component that
% follows it at the back of the packet.
% MC 4 and free-response (d) are the honest limit of the model: the empirical rule requires a
% mound-shaped, symmetric distribution, and the length-of-stay histogram is strongly skewed.
% Part (d) is also the spiral — it reaches back to Lesson 1.4 (shape, centre, spread) and 1.6
% (comparing distributions) for the language that decides whether a normal model is reasonable.
% GENERATED FROM ap_practice/main.tex. Each \writelines{n} in the blank is answered by an
% \ansline{} terse enough not to wrap past n lines — that keeps the page counts equal — and the
% \begin{work} block in 5(a) is carried over byte-identically.
% NO teachernote here — scoring guidance lives in the lesson plan.
\documentclass[10pt]{article}
\usepackage{apstats-article}
\usepackage{apstats-key}   % pulls in -boxes; do NOT also load -boxes

\begin{document}

\pageheader{Unit 1, Lesson 1.7}{AP Practice --- Answer Key}
% NAMESTRIP: no \namedateperiod here — mirrors the blank.

\vspace{0.15in}

% Authored BYTE-IDENTICALLY in ap_practice/ and ap_practice_key/ — this box is what keeps the
% blank and the key the same number of pages.
\boxguard[8]
\begin{remindbox}
\small
This page is \textbf{not required and not scored} --- your graded homework is the last section
of this packet. These are AP-format reps: work them for the practice, and bring what you get
stuck on to office hours or study hall.
\end{remindbox}

\vspace{0.12in}

\boxguard[26]
\begin{scenariobox}[Record A --- Newborn Birth Weights]{navy}
\small
At one hospital, the birth weights of newborns are approximately normal with mean
$\mu = 3{,}300$ grams and standard deviation $\sigma = 500$ grams. The curve below is marked at
the mean and at one, two, and three standard deviations either side of it.

\vspace{6pt}
\begin{center}
\begin{tikzpicture}[scale=0.9]
  \draw[navy, very thick] plot[smooth] coordinates
    {(-3.5,0.01) (-3,0.03) (-2.5,0.11) (-2,0.34) (-1.5,0.81) (-1,1.52) (-0.5,2.21) (0,2.50)
     (0.5,2.21) (1,1.52) (1.5,0.81) (2,0.34) (2.5,0.11) (3,0.03) (3.5,0.01)};
  \foreach \x/\h in {-3/0.03, -2/0.34, -1/1.52, 0/2.50, 1/1.52, 2/0.34, 3/0.03} {
    \draw[linegray, thin, dashed] (\x,0) -- (\x,\h);}
  \draw[->, thick] (-4.0,0) -- (4.0,0);
  \foreach \x/\lab in {-3/1800, -2/2300, -1/2800, 0/3300, 1/3800, 2/4300, 3/4800} {
    \draw (\x,-0.08) -- (\x,0.08); \node[below, font=\scriptsize] at (\x,-0.1) {\lab};}
  \node[font=\scriptsize] at (0,-0.62) {Birth weight (grams)};
\end{tikzpicture}
\end{center}
\end{scenariobox}

\vspace{0.12in}

\boxguard[26]
\begin{scenariobox}[Record B --- Length of Stay]{navy}
\small
The same hospital also recorded each newborn's \emph{length of stay}, in hours, for 300 births.
The mean stay is 46 hours and the standard deviation is 14 hours.

\vspace{6pt}
\begin{center}
\begin{tikzpicture}[scale=0.9]
  \foreach \i/\c in {0/12, 1/48, 2/76, 3/64, 4/45, 5/30, 6/20, 7/13, 8/8, 9/5, 10/3, 11/2} {
    \fill[goldacc] ({\i*0.5},0) rectangle ({(\i+1)*0.5},{\c*0.036});
    \draw[white, very thin] ({\i*0.5},0) rectangle ({(\i+1)*0.5},{\c*0.036});}
  \draw[->, thick] (0,0) -- (0,3.2);
  \draw[->, thick] (0,0) -- (6.4,0);
  \foreach \y/\lab in {0/0, 1.44/40, 2.88/80} {
    \draw (-0.08,\y) -- (0.08,\y); \node[left, font=\tiny] at (-0.1,\y) {\lab};}
  \foreach \x/\lab in {0/24, 1/36, 2/48, 3/60, 4/72, 5/84, 6/96} {
    \draw (\x,-0.08) -- (\x,0.08); \node[below, font=\scriptsize] at (\x,-0.1) {\lab};}
  \node[font=\scriptsize] at (3,-0.62) {Length of stay (hours)};
\end{tikzpicture}
\end{center}
\end{scenariobox}

\vspace{0.14in}

\boxguard[16]
\begin{headlinebox}{goldbg}
\small
\textbf{Section I --- Multiple Choice.} Circle the single best answer. Every answer choice is
about the context, not just the arithmetic --- read them all before choosing.
\end{headlinebox}

\vspace{0.10in}

\begin{enumerate}[leftmargin=1.9em, itemsep=0.13in]

  \item A newborn in Record A weighs 4,050 grams. What is this newborn's standardized score?
        \begin{enumerate}[label=\textbf{(\Alph*)}, leftmargin=2.2em, itemsep=2pt]
          \item $-1.5$
          \item $0.75$
          \item \ans{$1.5$}
          \item $2.5$
          \item $750$
        \end{enumerate}

  \item Using the empirical rule with Record A, approximately what percent of newborns weigh
        \textbf{less than} 2,300 grams?
        \begin{enumerate}[label=\textbf{(\Alph*)}, leftmargin=2.2em, itemsep=2pt]
          \item \ans{2.5\%}
          \item 5\%
          \item 16\%
          \item 32\%
          \item 95\%
        \end{enumerate}

  \item Newborn \emph{length} at the same hospital is approximately normal with mean 50.0 cm and
        standard deviation 2.0 cm. Baby A weighs 4,050 grams; Baby B is 52.0 cm long. Which baby
        is farther from average \emph{relative to its own distribution}?
        \begin{enumerate}[label=\textbf{(\Alph*)}, leftmargin=2.2em, itemsep=2pt]
          \item Baby A, because 4,050 is a far larger number than 52.0.
          \item \ans{Baby A, because its standardized score of 1.5 is larger than Baby B's 1.0.}
          \item Baby B, because its standardized score of 2.0 is larger than Baby A's 1.5.
          \item Baby B, because 52.0 cm is 2.0 cm above its mean, and 2.0 is larger than 1.5.
          \item Neither, because weight and length are measured in different units and cannot be
                compared at all.
        \end{enumerate}

  \item A nurse looks at Record B and claims that about 68\% of stays fall between 32 and 60
        hours, ``by the empirical rule.'' The best response is that
        \begin{enumerate}[label=\textbf{(\Alph*)}, leftmargin=2.2em, itemsep=2pt]
          \item the nurse is correct; the empirical rule applies to any distribution that has a
                mean and a standard deviation.
          \item the nurse is correct, because 32 and 60 are one standard deviation either side
                of the mean of 46.
          \item \ans{the empirical rule should not be used here, because the distribution of stays is
                strongly skewed to the right rather than mound-shaped and symmetric.}
          \item the empirical rule should not be used here, because 300 stays is too small a
                record to summarize.
          \item the empirical rule should not be used here, because length of stay was recorded
                in whole hours.
        \end{enumerate}

\end{enumerate}

\vspace{0.14in}

\boxguard[16]
\begin{headlinebox}{goldbg}
\small
\textbf{Section II --- Free Response.} Show your reasoning. Every answer must be written
\emph{in the context of the problem} --- that is what earns credit on the AP exam.
\end{headlinebox}

\vspace{0.10in}

\begin{enumerate}[leftmargin=1.9em, itemsep=0.16in, resume]

  \item Use Record A, and Record B in part (d).

        \vspace{6pt}
        \textbf{(a)} A newborn weighs 2,550 grams. Find this newborn's standardized score, then
        interpret it in one sentence in context.
        \begin{work}
        z &= \frac{x_i - \mu}{\sigma} \\[3pt]
          &= \frac{2550 - 3300}{500} = \frac{-750}{500} = -1.5
        \end{work}
        \ansline{This newborn weighs 1.5 standard deviations less than the mean birth}\\[1pt]
        \ansline{weight of 3,300 grams --- lighter than typical, but not extreme.}

        \vspace{6pt}
        \textbf{(b)} Use the empirical rule. Approximately what percent of newborns weigh between
        2,300 and 4,300 grams? Approximately what percent weigh more than 4,300 grams? Show how
        you got the second one.\\[4pt]
        \ansline{2,300 and 4,300 are two standard deviations either side of 3,300, so about}\\[1pt]
        \ansline{95\% of newborns weigh between them. That leaves about 5\% outside, and the}\\[1pt]
        \ansline{curve is symmetric, so about 2.5\% weigh more than 4,300 grams.}

        \vspace{6pt}
        \textbf{(c)} At a second hospital, birth weights are approximately normal with
        $\mu = 3{,}150$ grams and $\sigma = 400$ grams, and one newborn there weighs 3,750 grams.
        Compare that newborn's position within its own distribution to the 4,050 gram newborn's
        position within Record A.\\[4pt]
        \ansline{Second hospital: $z = (3750 - 3150)/400 = 1.5$. Record A: $z = 1.5$ as well.}\\[1pt]
        \ansline{The two newborns sit at exactly the same relative position --- each is 1.5}\\[1pt]
        \ansline{standard deviations above its own hospital's mean, despite different weights.}

        \vspace{6pt}
        \textbf{(d)} Look again at Record B. Is a normal model reasonable for length of stay?
        Name the features of the display that decide it, and say what your answer means for using
        the empirical rule on that variable.\\[4pt]
        \ansline{No. The histogram is unimodal but strongly skewed to the right, with a long}\\[1pt]
        \ansline{tail of stays past 72 hours and no matching tail on the left --- it is not}\\[1pt]
        \ansline{symmetric, so the empirical rule's 68--95--99.7 does not apply to these stays.}

\end{enumerate}

\vspace{0.12in}

\boxguard[12]
\begin{extensionbox}
\small
Standardized scores are how the SAT and the ACT are compared. Look up the mean and standard
deviation of each test's total score, then find the ACT score that sits at the same relative
position as an SAT score of 1300.\\[4pt]
\ansline{Answers vary with the year's figures. Method: standardize 1300 against the SAT}\\[1pt]
\ansline{mean and standard deviation, then undo it --- $x = \mu + z\sigma$ --- on the ACT scale.}
\end{extensionbox}

\vspace{0.12in}


\end{document}
